\pagebreak

\section{Xilinx netlist postprocessing module - xilinx/postprocess.3pl}\label{ch:pp}

    \index{library!xilinx/postprocess.3pl}
    This module converts the EDIF netlist produced by 3PL to a configuration
    file. It has a single procedure.
    
    This library must be incorporated by using a postmodule statement, e.g. -
    \begin{lstlisting}[gobble=4, language=threepl]
       postmodule ("xilinx/postprocess.3pl");
    \end{lstlisting}
    
    See \autorefph{sec:mpfcdecl} and \autorefph{sec:xpp}.
    
    More than one postprocessing module can be specified in which case an
    additional integer argument can be supplied to specify ordering. If omitted
    the order in which the modules have been added to the list will be the order
    of execution. Thus far only this module has been used.
    


    \subsection{library modules}

        None.



    \subsection{library procedures}

        \subsubsection{postprocess - post-process the netlist file}\label{sec:lppostprocess}

            {\bf Prototype:} \\
            \vsp
            \begin{lstlisting}[gobble=12, language=threepl]
                global procedure
                postprocess(
                    str(rname),
                    str(tool),
                    float(minver, 1.0),
                    float(maxver, 9999.0)
                ) {}
            \end{lstlisting}



            {\bf Output parameters:} \\
            
            None.

            {\bf Input parameters:} \\
            \begin{tabular}{| l | l | l | l |}
            \hline
                            & mode      & type   & \\
            \hline \hline
            {\bf rname}     & immediate & str    & optional file name for recipe file \\ \hline
            {\bf tool}      & immed     & str    & "ise" or "vivado" \\ \hline
            {\bf minver}    & immediate & float  & minimum version number  \\ \hline
            {\bf maxver}    & immediate & float  & maximum version number \\ \hline
            \end{tabular}

            {\bf In-line target execution:} \\
            no

            {\bf description:} \\
            This procedure executes the Xilinx place-and-route tools to
            generate the configuration bit file from the EDIF netlist file
            produced by 3PL. To do this it executes the Xilinx tools
            on the host CPU via calls to exec().
            
            {\bf rname} optional and is the name of a file containing a recipe
            of commands. If supplied it overrides the recipe file specified by
            directive {\bf postProcessRecipe}. Recipe files
            {\bf ise\_cpld.cmd}, {\bf ise\_fpga.cmd}, {\bf ise\_xflow.cmd}
            and {\bf vivado.tcl} are available in the directory
            containing postprocess.3pl.
            
            {\bf tool} "ise" or "vivado" must be supplied and specifies the
            main place-and-route tool to be used.
            
            If directives {\bf XILINX\_ISE} or {\bf XILINX\_VIVADO} are supplied
            (as appropriate to the platform) then this specifies the specific
            path of the place-and-route tool to execute. Otherwise the path
            specified by directive {\bf XILINX} will be used along with  {\bf
            minver} and {\bf maxver} to select the specific version of the ISE
            or Vivado software. The highest version \verb'>=' minver and
            \verb'<=' maxver will be used.


            if directive {\bf postProcessAppendFile} exists it specifies
            a file to append to the bit file, usually a .info file

            if directive {\bf postProcessCompress} exists and is true
            the bit file will be compressed.

            If directive {\bf postProcessFilter} exists it specifies that the
            standard output from the place-and-route tools should be filtered.
            Values may be -
            \blist
            \item   "error" - lines starting with "ERROR" are retained
            \item   "critical" - the above plus lines starting with "CRITICAL" are
                    retained
            \item   "warning" - the above plus lines starting with "WARNING" are
                    retained
            \item   "info" - the above plus lines starting with "INFO" are
                    retained
            \blend
            any other values being ignored. If this directive is not defined or
            has an invalid value then no filtering is applied.
    
    
    


    \subsection{library functions}
    
        None.


    \subsection{library classes}
    
        None.
